\part{Autopsie de formats représentatifs}
\section{Un format simple pour commencer : PPM}
Le format PPM binaire de Netpbm, variante P6, décrit une largeur,
une hauteur, une valeur maximale puis des composantes RGB~\cite{netpbm}.
Pour une image originale de deux pixels, rouge puis vert :
\begin{verbatim}
P6
2 1
255
[FF 00 00] [00 FF 00]
\end{verbatim}
Les crochets sont une annotation, pas des caractères à écrire.
L'en-tête est du texte, la charge utile est binaire. Une lecture
par lignes de texte appliquée aux pixels est donc incorrecte.
Avec ces seuls pixels, il n'existe ni miniature ni masque ni
profil ICC embarqué. Il faut connaître les conventions de
couleur et les limites de l'outil ; un tableau RGB brut est
beaucoup moins auto-descriptif qu'un format photographique riche.
La spécification PPM nominale décrit des composantes non linéaires
selon BT.709, pas des intensités linéaires. Des variantes rencontrées
en pratique utilisent pourtant sRGB ou du RGB linéaire : il ne faut
pas déduire une colorimétrie certaine de la seule extension
\texttt{.ppm}. Nos deux pixels extrêmes ne permettent pas de
distinguer ces fonctions de transfert.

Cet exemple permet d'isoler la différence entre matrice et
fichier : le tableau fait six octets, le fichier en fait
davantage à cause de son en-tête. Le script compagnon produit
ce fichier sans utiliser de bibliothèque d'image.

\section{TIFF : un graphe de répertoires et de données}
\subsection{Le contrat structurel}
TIFF classique commence par un ordre des octets, l'identifiant
42 et l'offset du premier IFD. Un IFD contient un nombre
d'entrées, des entrées de 12 octets et l'offset de l'IFD suivant.
Chaque entrée décrit un tag, un type, un nombre de valeurs et
une valeur courte ou un offset. Dimensions, profondeur,
compression, organisation et emplacement des bandes sont
des tags, pas un bloc universel de pixels placé à une adresse
constante~\cite{tiff}.

\subsection{En-tête original commenté octet par octet}
Un petit TIFF little-endian peut commencer ainsi :
\begin{verbatim}
Offset  Octets                  Interpretation
0       49 49                   "II", little-endian
2       2A 00                   entier 16 bits : 42
4       08 00 00 00             premier IFD a l'offset 8
8       09 00                   neuf entrees dans cet exemple
10      00 01 04 00 01 00 ...   debut du tag ImageWidth
\end{verbatim}
\texttt{00 01} vaut ici \texttt{0x0100}=256 : c'est
l'identifiant du tag ImageWidth, pas la largeur. Le type
\texttt{04 00} correspond à LONG ; le nombre de valeurs est
ensuite un entier sur quatre octets. Dans l'exemple du TP,
la valeur finale de cette entrée sera la largeur 2.

Le nombre 42 ne représente donc ni une taille de fichier ni
un nombre de pixels. Une entrée « largeur = 2 » doit être
distinguée du tag qui indique où chercher les pixels.

\subsection{Exemple de tags utiles}
\begin{longtable}{p{4.3cm}p{2cm}p{8.2cm}}
\toprule Nom & Identifiant & Rôle conceptuel \\ \midrule
ImageWidth / ImageLength & 256 / 257 & Dimensions de l'image concernée par cet IFD. \\
BitsPerSample & 258 & Profondeur par composante ; peut être une liste. \\
Compression & 259 & Méthode de reconstruction de la charge utile. \\
Photometric\-Interpretation & 262 & Sens des composantes : gris, RGB, palette, etc. \\
StripOffsets / StripByteCounts & 273 / 279 & Débuts et tailles des bandes. \\
Orientation & 274 & Transformation d'affichage. \\
SamplesPerPixel & 277 & Nombre de composantes à chaque position. \\
RowsPerStrip & 278 & Découpage vertical en bandes. \\
PlanarConfiguration & 284 & Composantes entrelacées ou séparées. \\
SampleFormat & 339 & Interprétation numérique : entier, flottant, etc. \\
\bottomrule
\end{longtable}
Le tableau sert de carte de lecture, pas de liste suffisante
pour rendre tout TIFF valide. Les valeurs par défaut, les
contraintes par mode et les extensions doivent être vérifiées
dans les documents normatifs.

\subsection{Bandes, tuiles, pages et sous-images}
Une image en bandes peut se lire par portions horizontales.
Une image en tuiles permet un accès spatial par blocs, utile
pour les très grandes images. Plusieurs IFD peuvent décrire
des pages ou des résolutions ; des SubIFD peuvent pointer vers
des images associées. L'ordre physique dans le fichier ne
définit pas forcément l'ordre logique des images.

Une compression sans perte conserve les échantillons stockés,
mais un TIFF peut aussi contenir une charge utile avec pertes.
Un TIFF flottant peut porter une image linéaire HDR ; un TIFF
entier 16 bits sRGB peut être une image SDR déjà rendue. Il
faut lire les tags et le profil au lieu de conclure à partir
de l'extension.

\subsection{BigTIFF et limite des offsets}
TIFF classique utilise de nombreux champs 32 bits, ce qui
limite l'adressage. BigTIFF introduit un en-tête et des
structures avec offsets plus larges~\cite{bigtiff}. Ce n'est
pas « un TIFF classique avec un plus gros fichier » : un
lecteur doit reconnaître la variante avant de lire ses IFD.
Un dépassement de capacité peut sinon produire des offsets
tronqués et un export apparemment réussi, mais inutilisable.

\section{PNG : chunks, prédiction et DEFLATE}
\subsection{Structure normative résumée}
PNG utilise une signature de huit octets puis des chunks.
Chaque chunk comporte longueur big-endian, type, données et
CRC sur type plus données. IHDR décrit notamment dimensions,
profondeur et type de couleur ; IDAT porte le flux zlib ; IEND
termine le fichier. Les IDAT concaténés forment un flux unique,
leurs frontières ne sont pas des frontières de lignes ou de
blocs DEFLATE. Les lignes décompressées portent un octet de
filtre à inverser. PNG prend en charge plusieurs types de
couleur et des profondeurs autorisées selon le type~\cite{png}.

\subsection{Un en-tête de deux pixels fabriqué pour le TP}
La signature en hexadécimal est :
\begin{verbatim}
89 50 4E 47 0D 0A 1A 0A
\end{verbatim}
Pour notre exemple RGB 8 bits de largeur 2 et hauteur 1,
les treize octets de données IHDR sont :
\begin{verbatim}
00 00 00 02   00 00 00 01   08 02 00 00 00
\end{verbatim}
La longueur du chunk vaut 13 ; son coût total est
$4+4+13+4=25$ octets. L'image de six octets utiles prend
donc nécessairement davantage de place que six octets dans
ce format. Le script calcule le CRC plutôt que d'en recopier
une constante opaque.

\subsection{Comprendre les filtres sans confondre octets et pixels}
Appelons $x$ l'octet courant, $a$ l'octet correspondant à gauche,
$b$ celui au-dessus et $c$ celui en haut à gauche. Un filtre
stocke $(x-P)\bmod256$ et le décodeur récupère
$(r+P)\bmod256$. Le voisin gauche est éloigné d'un pixel
complet en octets : pour RGB 8 bits, il est à trois octets,
pas forcément à l'octet précédent.

Exemple original : sur une ligne monochrome $(10,12,15)$,
la prédiction par le voisin gauche donne $(10,2,3)$.
Le premier voisin absent vaut zéro dans ce modèle. Une
prévision par la ligne précédente est efficace sur des
zones verticalement régulières. Une prédiction combinant
gauche, dessus et diagonale peut mieux suivre les gradients.
Le filtre ne supprime pas les différences : il les réexprime.

Les cinq types usuels sont None, Sub, Up, Average et Paeth.
Le TP implémente leur inversion sur des octets et vérifie
leur réversibilité sur des données synthétiques. Il ne
prend volontairement pas en charge toute la norme.

\subsection{Sans perte ne signifie pas identique en octets}
Deux encodeurs peuvent choisir des filtres, niveaux de
compression et découpages de chunks différents. Les fichiers
auront des SHA-256 distincts tout en redonnant le même RGB.
Une suppression du profil ou une conversion de couleur avant
l'encodage peut modifier le rendu, même si le codec PNG reste
sans perte sur les pixels qu'il a reçus.

\section{JPEG : marqueurs et flux de coefficients}
\subsection{Fichier et procédé de compression}
JPEG est une famille de procédés, pas exclusivement la variante
baseline avec pertes la plus courante. Les fichiers usuels
associent ces procédés à des conventions comme JFIF ou Exif.
On rencontre notamment SOI, DQT, SOF, DHT, SOS et EOI pour
l'entrée, les tables, les caractéristiques, les scans et la fin.
Les modes progressifs transmettent les coefficients en
plusieurs scans ; des procédés prédictifs sans perte existent
aussi dans T.81~\cite{jpeg}.

\subsection{Schéma de lecture d'un JPEG usuel}
\begin{verbatim}
SOI
  APP... : description JFIF, Exif, ICC, XMP selon conventions
  DQT    : tables de quantification
  SOF... : dimensions, composantes, facteurs d'echantillonnage
  DHT    : tables Huffman si ce mode les utilise
  SOS    : description du scan
  ...    : donnees entropiques du scan
  ...    : scans ou segments supplementaires selon le mode
EOI
\end{verbatim}
Ce schéma ne fixe pas un ordre unique pour tous les modes.
Dans les données entropiques, des conventions d'échappement
et de redémarrage évitent de prendre des octets de données
pour des marqueurs. Chercher simplement \texttt{FF D9}
dans n'importe quelle sous-structure n'est pas un parseur.

Une miniature JPEG peut elle-même être contenue dans un
JPEG ou un RAW plus grand. Il faut distinguer le flux principal
de l'aperçu. Un logiciel peut réussir à afficher l'aperçu
tout en échouant à décoder la charge utile principale.

\section{DNG : TIFF plus une description du négatif numérique}
DNG étend TIFF avec des informations de calibration et de
développement : CFA, zones actives, niveaux, matrices et
autres tags de pipeline. La spécification distingue notamment
des données mosaïquées et des données linéaires déjà
dématricées ; plusieurs compressions sont possibles selon
les versions et modes. Les tags d'un profil DCP décrivent
une transformation ou un rendu, pas les pixels de la
photographie~\cite{dng}.

Un DNG valide doit être cohérent comme conteneur \emph{et}
comme négatif numérique. Écrire trois canaux RGB dans un TIFF
puis ajouter une extension DNG ne recrée pas la mosaïque
originale ni toutes ses possibilités de développement.
Un DNG linéaire peut être légitime, mais il doit déclarer
correctement sa nature et ses calibrations. Si des effets
sont déjà incorporés aux pixels, il ne faut pas décrire
ces pixels comme s'ils étaient encore les mesures brutes.

\section{HEIF, HEIC, AVIF et conteneurs à boîtes}
HEIF organise des éléments d'image et leurs relations dans
une structure fondée sur ISO BMFF. HEIC est couramment
associé à des images codées en HEVC ; AVIF associe une image
AV1 à une organisation HEIF. La spécification AVIF décrit
notamment propriétés de couleur et relations d'éléments~\cite{avif}.

Une boîte peut commencer par une taille et un type de quatre
caractères, puis contenir des sous-boîtes ou une charge utile.
Des associations peuvent relier une image à une miniature,
un alpha ou d'autres données. Un parser doit savoir quelle
structure est un conteneur avant de descendre récursivement :
tous les octets d'une boîte ne sont pas d'autres boîtes.

L'existence de modes HDR, alpha ou haute précision ne signifie
pas que tout lecteur les gère. Vérifier les propriétés exactes
du fichier et la compatibilité de la chaîne de décodage.
Un codec moderne n'est pas automatiquement un espace couleur.

\section{OpenEXR et images de calcul scientifique ou de rendu}
OpenEXR permet des canaux nommés, notamment HALF, FLOAT et
UINT, des fenêtres de données, tuiles et parties multiples.
Certaines compressions sont sans perte, d'autres avec
pertes. Des données profondes peuvent contenir plusieurs
échantillons par position ; les composantes couleur sont
habituellement linéaires et prémultipliées dans les usages
de compositing~\cite{openexr}.

Un canal Z peut représenter la profondeur et non une couleur.
Une composante supérieure à 1 peut porter une énergie
relative élevée sans être « invalide ». Elle ne devient
blanche sur un écran qu'après choix d'un rendu. L'architecture
d'un fichier image dépasse ainsi largement le triplet RGB.

\section{Panorama et choix selon l'usage}
\begin{longtable}{p{2.1cm}p{4.1cm}p{4.4cm}p{3.9cm}}
\toprule Famille & Atout fréquent & Attention & Usage typique \\ \midrule
JPEG usuel & Taille modeste, diffusion & Pertes et chroma possibles & Livraison de photos SDR. \\
PNG & Réversibilité des échantillons & Coût sur photographies bruitées & Graphiques, captures, alpha. \\
TIFF & Souplesse et haute précision & Variantes nombreuses & Échange et comparaison contrôlée. \\
RAW constructeur & Mesures et calibration & Support dépendant du modèle/mode & Développement photographique. \\
DNG & Description documentée du négatif & Cohérence et version à respecter & Échange/archivage RAW ou linéaire. \\
AVIF/\allowbreak HEIF & Codage moderne, éléments associés & Compatibilité variable & Diffusion et usages HDR compatibles. \\
OpenEXR & Flottants et canaux génériques & Pas un rendu photo universel & VFX, intermédiaires, calcul HDR. \\
\bottomrule
\end{longtable}
Il existe aussi BMP, GIF, WebP, JPEG~2000, JPEG~XL, PSD et
des formats médicaux ou géographiques. Leurs objectifs
diffèrent : animation, couches, séries, coordonnées,
profondeur ou accès aléatoire. Les principes étudiés ici
restent utiles pour les aborder sans supposer qu'ils suivent
tous le modèle JPEG ou TIFF.
